
\uline{le sous-schéma fermé de} @@M0@@ \uline{défini par} @@M1@@, @@M2@@ \uline{le schéma obtenu en faisant éclater} @@M3@@, @@M4@@ \uline{le fibré projectif associé à} @@M5@@, @@M6@@ \uline{le morphisme structural} @@M7@@. \uline{Sur} @@M8@@, \uline{on considère la suite exacte canonique} :
@@M9@@
\uline{et le morphisme composé} :
@@M10@@

\uline{Alors} :
\begin{enumerate}[label=\roman*)]
\item \uline{La formation de} @@M11@@ \uline{commute à tout changement de base} @@M12@@ \uline{tel que} @@M13@@ \uline{soit régulier}.
\item \uline{L'immersion fermée} @@M14@@ \uline{définie par l'homomorphisme surjectif} @@M15@@ \uline{défini par} @@M16@@ \uline{en degré} @@M17@@ \uline{est une immersion régulière}, \uline{son idéal étant l'image du morphisme régulier} @@M18@@ ; \uline{en particulier si} @@M19@@ \uline{est constante}, @@M20@@.
\item \uline{Le} @@M21@@-\uline{Module} @@M22@@ \uline{est parfait en tant que} @@M23@@-\uline{Module gradué}, \uline{i}.\uline{e}. \uline{il existe localement une résolution finie}, \uline{graduée}, \uline{de} @@M24@@ \uline{par des} @@M25@@-\uline{Modules gradués localement libres de type fini} (VI 1.2).
\end{enumerate}

i) Comme @@M26@@, le résultat est une conséquence immédiate des relations (1.2.3) et de EGA II.

ii) Les propriétés énoncées étant locales sur @@M27@@, on peut supposer @@M28@@ affine, et de plus @@M29@@ libre sur @@M30@@. On pose @@M31@@.

On fait d'abord la démonstration dans le cas standard @@M32@@, @@M33@@ étant muni d'une base @@M34@@, avec @@M35@@. Alors @@M36@@, la graduation étant prise par rapport aux @@M37@@. Pour montrer que @@M38@@ est un morphisme régulier dont l'image est l'Idéal de @@M39@@, il suffit de le faire localement sur @@M40@@, donc sur les
